\documentclass[11pt,a4paper]{article}
\newcommand{\doctag}{Toolchain migration journal}
\usepackage[a4paper,margin=2.25cm,headheight=15pt]{geometry}
\usepackage[T1]{fontenc}
\usepackage{lmodern,microtype,graphicx,booktabs,array,tabularx,amsmath}
\usepackage[table]{xcolor}
\usepackage{enumitem,listings,fancyhdr}
\usepackage[hidelinks]{hyperref}
\definecolor{navy}{HTML}{163A5F}
\definecolor{pale}{HTML}{EDF3F8}
\definecolor{muted}{HTML}{526272}
\setlength{\parindent}{0pt}
\setlength{\parskip}{0.48em}
\setlength{\emergencystretch}{2em}
\renewcommand{\arraystretch}{1.16}
\newcolumntype{Y}{>{\raggedright\arraybackslash}X}
\newcolumntype{P}[1]{>{\raggedright\arraybackslash}p{#1}}
\setlist{nosep,leftmargin=1.5em,topsep=0.35em}
\pagestyle{fancy}
\fancyhf{}
\fancyhead[L]{\small ERTS Assignment 2\quad\textcolor{muted}{\doctag}}
\fancyhead[R]{\small 11 October 2026}
\fancyfoot[R]{\small\thepage}
\renewcommand{\headrulewidth}{0.4pt}
\lstset{basicstyle=\ttfamily\footnotesize,breaklines=true,columns=fullflexible,
  keepspaces=true,showstringspaces=false,frame=single,rulecolor=\color{muted!30},
  backgroundcolor=\color{pale},xleftmargin=0.2em,xrightmargin=0.2em}
\newcommand{\banner}[1]{\colorbox{pale}{\parbox{\dimexpr\linewidth-2\fboxsep\relax}{#1}}}
\newcommand{\docopening}[2]{
 {\Large\bfseries\color{navy} Embedded Real Time Systems\par}
 \vspace{0.35em}{\LARGE\bfseries #1\par}
 \vspace{0.4em}{\large #2\par}
 \vspace{0.5em}Firas Harbo Saleh \quad Student ID: 202109174\par
 \textcolor{muted}{Review snapshot: \texttt{assignment\_2/Firas}, commit \texttt{c38a216}.}\par
 \vspace{0.6em}}
\newcommand{\src}[1]{{\small\path{#1}}}

\begin{document}
\docopening{Toolchain Migration Journal}{Course Vivado 2017.2 / SDK to native Vivado and Vitis 2026.1}

\banner{\textbf{Purpose.} Record the differences that affect this assignment,
the post-review changes, why they were made, and remaining acceptance boundaries.
This is an observed project migration, not a benchmark
of all Vivado releases or a substitute for course approval.}

\section*{1. Starting point and acceptance boundary}
The course assignment explicitly requires Vivado 2017.2 (p.~2), and its
installation guide includes Software Development Kit (SDK). The supplied
Ex7 project also records 2017.2 and the original ZYBO part
\texttt{xc7z010clg400-1}. Our implementation uses installed native
Vivado/Vitis 2026.1 on Ubuntu 26.04 LTS, targeting that same part.

The installation guide was tested on Windows 10 but offers a Windows or Linux
installer. Linux is therefore not itself a violation of the written guide;
using 2026.1 instead of 2017.2 is the explicit version deviation. A successful
board run does not establish instructor acceptance of that deviation.

\begin{tabularx}{\linewidth}{P{3.15cm} Y Y}
\toprule Layer & Course / original project & Current implementation \\
\midrule
Hardware design & Vivado 2017.2 project and board preset & Vivado 2026.1 project/IP upgrade, explicit constraints \\
HLS & Vivado HLS, guided GUI steps & Vitis HLS, Tcl-controlled Vivado IP flow \\
Hardware handoff & SDK export using HDF & Fixed XSA including bitstream \\
Embedded application & SDK workspace and legacy BSP & Vitis Unified platform, standalone domain, generated SDT BSP \\
Automation & SDK-era project/debug workflow & Vitis Python build plus XSDB Tcl board launch \\
Peripheral lookup & Legacy device-ID configuration & Base-address/compatible-name configuration where adapted \\
ANN register interface & One original AXI-Lite slave & Two generated slaves: control and array storage \\
Physical target & Original ZYBO XC7Z010 & Same ZYBO; no Z7 board substitution \\
\bottomrule
\end{tabularx}

\section*{Evidence conventions}
\textbf{Observed} changes have local source/artifact evidence.
\textbf{Toolchain changes} enable the new workflow; \textbf{independent
corrections} address algorithms or robustness. The reviewed commit is the
starting point; this revision includes subsequent local source fixes.

No 2017.2 toolchain was rerun. Historical comparisons use the course material
and supplied artifacts, not a new performance comparison between releases.

\newpage
\section*{2. HLS and hardware-project changes}
\subsection*{Scripted HLS with a controlled baseline}
\texttt{run\_hls.tcl} creates baseline and pipelined solutions for the same
XC7Z010 part and 10 ns target. Both explicitly use \texttt{-flow\_target vivado}
so the output is reusable IP, rather than a Vitis acceleration kernel flow.
Both set \texttt{config\_compile -pipeline\_loops 0}; only the pipelined
solution adds II=1 directives to the two named inner dot-product loops.

AMD's migration documentation records changes in automatic loop-pipelining and
clock-uncertainty defaults between historical Vivado HLS and Vitis HLS [W3].
Explicit settings make our two-solution comparison interpretable; a tool
upgrade alone must not be mistaken for the effect of a deliberate pipeline.
The current reports show 2.70 ns clock uncertainty. This is an observed 2026.1
report value, not a reconstructed 2017.2 setting.

The flow performs C simulation and synthesis for both solutions, VHDL
co-simulation for the lower-latency solution, and IP-catalog export with
\texttt{-rtl vhdl -version 1.0.1}. The export log and settings retain the
requested version; IP XML represents it as version 1.0 plus core revision 1.
Reading the XML's version field alone would incorrectly suggest the request
was ignored.

\subsection*{Project, IP and directory migration}
The original Ex7 XPR names the board preset
\texttt{digilentinc.com:zybo:part0:1.0}. The migrated XPR keeps the silicon part
but no longer depends on that unavailable local preset. The block design
changes from the original XML representation to a 2026.1 JSON representation,
and the processing-system, GPIO, interconnect and reset IP revisions change.

\texttt{build\_ex7\_2026.tcl} selects the newly exported HLS-IP repository,
sets the button GPIO to a custom four-bit input interface, attempts the IP
upgrades, connects the switch/LED ports and new array bus, assigns addresses,
validates the BD, generates targets and runs implementation through bitstream
generation. It writes timing/utilization reports and a bitstream-containing XSA.

The script uses the constraints already stored in the project; it does not
generate those pin assignments itself. The change from board automation to
explicit ports/pins is a practical repair for this machine's missing old board
preset, not a rule that every 2026.1 project must avoid board files.

\subsection*{Migration is not merely opening an old project}
Accepting an upgrade dialog updates project metadata, but does not prove IP
compatibility, correct pins, usable BSP drivers or application correctness.
For this project the integration repairs, new BSP, driver configuration and
physical UART check were separate necessary steps. Directory layout changes
and renamed generated files also mean that old paths in SDK-era instructions
cannot simply be copied unchanged.

For Exercises 3--5, a direct upgrade encountered locked legacy board-preset
outputs. The new \texttt{build\_ex345\_2026.tcl} reconstructs the supplied Ex5
block design from its original Tcl using current catalog revisions, unchanged
custom LED/matrix RTL, and explicit ZYBO pins. This Ex3-style processor/IO
platform with added matrix IP hosts three separate applications.

\newpage
\section*{3. Software platform, drivers and board connections}
\subsection*{HDF / SDK to XSA / Vitis Unified}
The original retained Vivado log copies a system definition to HDF and launches
SDK. The new hardware script uses \texttt{write\_hw\_platform -fixed -include\_bit}
to produce \texttt{Ex7\_2026\_1.xsa}. The native Python builder creates a
standalone platform/domain for \texttt{ps7\_cortexa9\_0}, imports the application,
shared timer header and AMD platform support, and builds the ELF.

The important distinction is between project automation and debugging.
AMD states that XSCT is no longer available starting with 2026.1 and recommends
the Python APIs for project management; XSDB remains available for debug [W1].
Accordingly our project build uses Vitis Python, while our working board launch
uses Tcl XSDB. These are complementary tools, not contradictory workflows.

\subsection*{SDT lookup semantics}
The newer system-device-tree (SDT) configuration no longer identifies devices
with the old \texttt{DeviceId} field; compatible names and base addresses are
used [W2]. The shared timer selects the base-address lookup under \texttt{SDT}
and retains the old device-ID path otherwise. Ex3 now applies this distinction
to the switch GPIO. Separately, custom LED-IP access uses register 0, not an
AXI GPIO driver. Semantic generated macros select peripherals without relying
on arbitrary numbered GPIO instances. Ex345 LEDs are at \texttt{0x43C10000},
unlike the original Ex3-only \texttt{0x43C00000} mapping.

\subsection*{Two ANN AXI-Lite buses and a project-specific workaround}
The supplied old IP has one AXI-Lite interface. Our current generated IP splits
control and array storage, despite the shared source bundle name. The saved
BD and generated driver establish the following mapping:

\begin{center}
\begin{tabular}{lll}
\toprule Interface & Function & Base address \\
\midrule
\texttt{s\_axi\_BUS} & Start/done/control & \texttt{0x43C00000} \\
\texttt{s\_axi\_BUS\_r} & Input, weights and output storage & \texttt{0x40000000} \\
\texttt{buttons/S\_AXI} & AXI GPIO & \texttt{0x41200000} \\
\bottomrule
\end{tabular}
\end{center}
The generated SDT config table lists the ANN addresses in ascending order,
but the driver fields expect control first and arrays second. The application's
\texttt{initializeAccelerator} constructs the named fields explicitly:
\texttt{Bus\_BaseAddress = BASEADDR\_1} and
\texttt{Bus\_r\_BaseAddress = BASEADDR}. This is verified for this generated
artifact. New compile-time/runtime guards check these project-specific addresses
and verify that driver memory widths/depths match the shared dimension header.

\subsection*{Physical pins stay tied to the actual board}
LED0--3 remain M14/M15/G14/D18; SW0--3 remain G15/P15/W13/T16.
The XDC also activates existing button pins under the generated port name.
The generated VHDL continuously connects switches to LEDs. Correct constraints
and numerical inference output do not replace the required physical LED test.

\newpage
\section*{4. Changes that are not caused by the tool version}
These changes should be credited as implementation work, not presented as
unavoidable Vivado migration requirements.

\begin{tabularx}{\linewidth}{P{3.2cm} Y}
\toprule Change & Reason and consequence \\
\midrule
Weight row strides & Uses \texttt{INPUT\_LEN*i} for input weights and \texttt{NEURONS\_LEN*i} for output weights. These are the contiguous fully connected layouts, unlike the overlapping rows in the starter reference. \\
48-bit accumulators & Protect fixed-point dot-product accumulation before shifts, particularly the second layer. This is arithmetic design, not an IDE compatibility patch. \\
Extended testbench & Preserves the original; adds an explicitly adapted course vector and four signed/zero/overflow cases with an independent integer reference. \\
Aligned transfer buffers & Packs byte/short arrays into word-aligned buffers before driver calls. Avoids relying on casts from differently aligned arrays. \\
Transfer and polling checks & Checks word counts and bounds ready/done polling. Makes failures visible rather than permitting an indefinite wait. \\
Separated measurements & Retains CPU/load/run timers; adds full packing-to-readback latency and ten-trial min/mean/max at application \texttt{-O2}. \\
\bottomrule
\end{tabularx}

\subsection*{The reference discrepancy must be visible}
A direct host build of the untouched supplied testbench against the current
core exits 1. Output 0 is 331 versus 41.003418 in that reference; all ten outputs
disagree. The current custom test passes, and a temporary 8/3/4-dimension host
variant also passes. These facts can coexist because the tests describe
different weight layouts and arithmetic references.

The implemented remedy is a documented course-derived adapted reference,
preserving the original and identifying changed strides and fixed-point
expectations. Default and 8/3/4 host suites pass; the untouched reference still
fails. Acceptance of the adaptation remains a course decision.
Neither silently replacing the reference nor changing the core to overlapping
weight rows merely to satisfy that reference is a defensible migration story.

\subsection*{Historical pre-fix performance, with the correct interpretation}
Baseline/pipelined HLS latency is 33,056/6,612 cycles; both are 2026.1 builds.
VHDL co-simulation measures 6,610 cycles. Pipelining increases DSP from 2 to 3,
FF from 638 to 1,112 and LUT from 716 to 1,126; BRAM18K stays at 10.

The fresh board run reports CPU 186,747 ticks, loading 233,849 ticks and
start-to-done 21,700 ticks. At 325 MHz these are 574.6, 719.5 and 66.8 $\mu$s.
The $8.61\times$ compute-only comparison is against an \texttt{-O0 -g3} ARM
application, not optimized CPU performance. The load-plus-run subtotal is
255,549 ticks; packing, setup and output readback are outside it. Loading alone
already exceeds software compute, demonstrating the one-shot transfer bottleneck.

\newpage
\section*{5. Reproducing the native workflow safely}
The commands below document the revised native flow. Run each stage sequentially, with
only one hardware/build session using the project at a time.

\textbf{Before rebuilding:} preserve canonical sources outside generated folders.
HLS anchors its working directory and renames the old project as a backup;
Vivado resets synthesis. Ex7 Vitis reuses its workspace with source backups, or
renames it on explicit \texttt{--rebuild}. Ex345 also preserves old workspaces.

\subsection*{Stage A: regenerate HLS IP}
From the repository root \texttt{ERTS\_assigments}:
\begin{lstlisting}
source /home/firas/AMD/2026.1/Vitis/settings64.sh
vitis-run --mode hls --tcl \
  Assignment_2/exercise7_neural_net/hls/run_hls.tcl
\end{lstlisting}
Check both C simulations, synthesis tables and VHDL co-simulation. Confirm the
generated IP exists before updating the hardware project.

\subsection*{Stage B: implement and export hardware}
From the same repository root, using a separate shell for Vivado settings:
\begin{lstlisting}
source /home/firas/AMD/2026.1/Vivado/settings64.sh
vivado -mode batch -source \
  Assignment_2/hardware/build_ex7_2026.tcl
\end{lstlisting}
Check BD validation, implementation completion, timing and bitstream/XSA
creation. This is a batch process, not a reason to open another GUI. The current
scripts cap implementation jobs and tool threads at two.

\subsection*{Stage C: build the standalone application}
From the repository root in the Vitis-configured shell:
\begin{lstlisting}
vitis -s Assignment_2/exercise7_neural_net/sdk/build_vitis_2026.py
\end{lstlisting}
Check the generated address map, ANN config fields and compiler flags, then
confirm \src{Assignment_2/hardware/Ex7_vitis/neuralNetTest/build/neuralNetTest.elf}.

\subsection*{Stage D: run on the identified ZYBO}
Reuse an existing \texttt{hw\_server} on port 3121. Identify the expected cable
and XC7Z010 first, then capture UART at 115200 baud and run:
\begin{lstlisting}
python3 Assignment_2/hardware/capture_board.py 7 \
  --output ../outputs/fix_evidence/board_ex7_new_run
\end{lstlisting}
This resets and programs volatile board state, initializes the PS, downloads
the ELF and starts the CPU; it does not write flash. The launcher selects the
known cable and rejects absent/ambiguous targets. Real UART bytes, timestamped
events, XSDB output and artifact hashes are retained, not reconstructed images.

\newpage
\section*{6. Provenance, remaining work and sources}
\small
\begin{tabularx}{\linewidth}{P{2.3cm} Y}
\toprule Commit & Observed role \\
\midrule
\texttt{47e3b21} & Initial exercise implementation, fixed-point core, tests and software logic. \\
\texttt{29f6373} & Main 2026.1 Ex7 hardware-project migration and build flow. \\
\texttt{cbec41c} & Native Vitis build, SDT compatibility, dual-bus address handling, transfer hardening and board launch. \\
\texttt{c38a216} & Recorded successful physical benchmark; reviewed repository snapshot. \\
\bottomrule
\end{tabularx}
\normalsize

\subsection*{Changes implemented after review}
\begin{enumerate}
\item Correct custom LED-IP access and SDT GPIO lookup; explicitly select an
OCM linker script and verify the built ARM ELF heap/stack locations.
\item Restore prescribed matrix types/signatures and literal transpose inputs,
retaining the previous fixture as a regression case.
\item Preserve the original neural reference and add a documented adapted
course vector plus signed, zero-state and wide-accumulation tests.
\item Build applications at \texttt{-O2}; report ten ANN trials and full
transaction timing alongside the separate loading and compute measurements.
\item Share dimensions across C/C++, round transfer buffers up to whole words,
zero padding, and reject stale generated driver dimensions at compile time.
\item Preserve workspaces, guard addresses, select the cable serial and retain
artifact hashes with raw UART evidence; update stale build instructions.
\end{enumerate}
These are local post-review fixes. The user confirmed the full physical LED
counter during the 11 October rerun. Separate switch mapping observations and course
acceptance of 2026.1 and the adapted reference remain outstanding. The hand-in
and review closure addendum record the latest completed verification runs.

\subsection*{Local evidence}
Course PDFs: \src{work/assignment2/Assignment_2/ERTS_Assignment2_ANN_ZYBO.pdf}
and \src{/home/firas/Downloads/ERTS_Vivado_Installation_2017_2.pdf}.
Original project: \src{/home/firas/Downloads/Ex7.xpr.zip}; SHA-256:
{\footnotesize\path{6928f168497e39fdb69fe8ebe2b08604eae389b49132aec17fc6949fabaf4cac}}.

Under \src{ERTS_assigments/Assignment_2/}, primary implementation evidence is
\src{hardware/build_ex7_2026.tcl}, \src{hardware/Ex7/Ex7.xpr}, its saved
\texttt{system.bd}/XDC, \src{exercise7_neural_net/hls/run_hls.tcl},
\src{exercise7_neural_net/sdk/build_vitis_2026.py},
\src{exercise7_neural_net/sdk/neuralNetTest.c}, and
\src{common/zynq_timer.h}. Generated \texttt{xneuralnet\_g.c}, header and
\texttt{xparameters.h} substantiate the bus-field workaround.
Historical physical evidence is under
\src{outputs/native_screenshots/zybo_retry_20261011T085514Z/}; its blank screenshot
attempt is excluded. Post-review build and UART evidence is under
\src{outputs/fix_evidence/}, including the recorded user counter observation.

\subsection*{AMD primary references (accessed 11 October 2026)}
\small
[W1] \href{https://docs.amd.com/r/en-US/ug1400-vitis-embedded/XSCT-to-Python-API-Migration}{UG1400, 2026.1: XSCT to Python API Migration}.
Supports project-automation changes and continued XSDB debugging.\par
[W2] \href{https://docs.amd.com/r/2024.2-English/ug1647-porting-embeddedsw-components/What-s-New}{UG1647, 2024.2: What's New in the SDT-based flow}.
Supports device-ID versus compatible-name/base-address semantics.\par
[W3] \href{https://docs.amd.com/r/2021.1-English/ug1399-vitis-hls/Default-User-Control-Settings}{UG1399, 2021.1: Default User Control Settings}.
Historical migration context for HLS defaults; current numerical settings
are taken from this project's 2026.1 reports, not inferred from this older page.
\end{document}
